% Options for packages loaded elsewhere
\PassOptionsToPackage{unicode}{hyperref}
\PassOptionsToPackage{hyphens}{url}
\PassOptionsToPackage{dvipsnames,svgnames,x11names}{xcolor}
\documentclass[
  11pt,
]{article}
\usepackage{xcolor}
\usepackage[margin=1in]{geometry}
\usepackage{amsmath,amssymb}
\setcounter{secnumdepth}{-\maxdimen} % remove section numbering
\usepackage{iftex}
\ifPDFTeX
  \usepackage[T1]{fontenc}
  \usepackage[utf8]{inputenc}
  \usepackage{textcomp} % provide euro and other symbols
\else % if luatex or xetex
  \usepackage{unicode-math} % this also loads fontspec
  \defaultfontfeatures{Scale=MatchLowercase}
  \defaultfontfeatures[\rmfamily]{Ligatures=TeX,Scale=1}
\fi
\usepackage{lmodern}
\ifPDFTeX\else
  % xetex/luatex font selection
\fi
% Use upquote if available, for straight quotes in verbatim environments
\IfFileExists{upquote.sty}{\usepackage{upquote}}{}
\IfFileExists{microtype.sty}{% use microtype if available
  \usepackage[]{microtype}
  \UseMicrotypeSet[protrusion]{basicmath} % disable protrusion for tt fonts
}{}
\makeatletter
\@ifundefined{KOMAClassName}{% if non-KOMA class
  \IfFileExists{parskip.sty}{%
    \usepackage{parskip}
  }{% else
    \setlength{\parindent}{0pt}
    \setlength{\parskip}{6pt plus 2pt minus 1pt}}
}{% if KOMA class
  \KOMAoptions{parskip=half}}
\makeatother
\usepackage{longtable,booktabs,array}
\newcounter{none} % for unnumbered tables
\usepackage{calc} % for calculating minipage widths
% Correct order of tables after \paragraph or \subparagraph
\usepackage{etoolbox}
\makeatletter
\patchcmd\longtable{\par}{\if@noskipsec\mbox{}\fi\par}{}{}
\makeatother
% Allow footnotes in longtable head/foot
\IfFileExists{footnotehyper.sty}{\usepackage{footnotehyper}}{\usepackage{footnote}}
\makesavenoteenv{longtable}
\setlength{\emergencystretch}{3em} % prevent overfull lines
\providecommand{\tightlist}{%
  \setlength{\itemsep}{0pt}\setlength{\parskip}{0pt}}
\usepackage{bookmark}
\IfFileExists{xurl.sty}{\usepackage{xurl}}{} % add URL line breaks if available
\urlstyle{same}
\hypersetup{
  pdftitle={An exact right-angled unitarity triangle with λ = 1/√20 and A = √(2/3): a testable parametrisation of quark mixing},
  pdfauthor={Matthew Gelfman (ORCID 0009-0000-6402-2258); Melissa Behrends-Gelfman (ORCID 0009-0001-6281-1947)},
  colorlinks=true,
  linkcolor={Maroon},
  filecolor={Maroon},
  citecolor={Blue},
  urlcolor={Blue},
  pdfcreator={LaTeX via pandoc}}

\usepackage{amssymb}
\usepackage{newunicodechar}
\newunicodechar{°}{\ensuremath{^\circ}}
\newunicodechar{±}{\ensuremath{\pm}}
\newunicodechar{²}{\ensuremath{^{2}}}
\newunicodechar{³}{\ensuremath{^{3}}}
\newunicodechar{¹}{\ensuremath{^{1}}}
\newunicodechar{×}{\ensuremath{\times}}
\newunicodechar{α}{\ensuremath{\alpha}}
\newunicodechar{β}{\ensuremath{\beta}}
\newunicodechar{γ}{\ensuremath{\gamma}}
\newunicodechar{δ}{\ensuremath{\delta}}
\newunicodechar{λ}{\ensuremath{\lambda}}
\newunicodechar{μ}{\ensuremath{\mu}}
\newunicodechar{σ}{\ensuremath{\sigma}}
\newunicodechar{χ}{\ensuremath{\chi}}
\newunicodechar{⁵}{\ensuremath{^{5}}}
\newunicodechar{⁶}{\ensuremath{^{6}}}
\newunicodechar{⁻}{\ensuremath{^{-}}}
\newunicodechar{₁}{\ensuremath{_{1}}}
\newunicodechar{₂}{\ensuremath{_{2}}}
\newunicodechar{₃}{\ensuremath{_{3}}}
\newunicodechar{−}{\ensuremath{-}}
\newunicodechar{√}{\ensuremath{\surd}}
\title{An exact right-angled unitarity triangle with \ensuremath{\lambda} = 1/\ensuremath{\surd}20 and A =
\ensuremath{\surd}(2/3): a testable parametrisation of quark mixing}
\author{Matthew Gelfman (ORCID 0009-0000-6402-2258) \and Melissa
Behrends-Gelfman (ORCID 0009-0001-6281-1947)}
\date{25 November 2026}

\begin{document}
\maketitle

\subsection{Abstract}

We take the right angle of the unitarity triangle at \ensuremath{\alpha} as exact, as
quark-mixing sum rules allow (Antusch, King, Malinsky and Spinrath
2010), and choose simple values for the other three quark-mixing
parameters: \ensuremath{\lambda} = 1/\ensuremath{\surd}20, A = \ensuremath{\surd}(2/3), \ensuremath{\beta} = 22.5\ensuremath{^\circ} and \ensuremath{\gamma} = 67.5\ensuremath{^\circ}. The result
is an exactly unitary CKM matrix (Kobayashi and Maskawa 1973) with
Jarlskog invariant J = 2.94 \ensuremath{\times} 10\ensuremath{^{-}}\ensuremath{^{5}} (Jarlskog 1985) and sin 2\ensuremath{\beta} = 1/\ensuremath{\surd}2.
All seven directly measured quantities lie within 1.6\ensuremath{\sigma}, with pulls from
\ensuremath{-}1.2\ensuremath{\sigma} (\textbar V\textsubscript{ub}\textbar) to +1.6\ensuremath{\sigma} (\ensuremath{\alpha}), giving \ensuremath{\chi}\ensuremath{^{2}} =
6.0 for seven quantities. The four values were chosen, not derived, so
the \ensuremath{\chi}\ensuremath{^{2}} is not evidence. We present the form as a parametrisation with
sharp, near-term tests. The matrix is exactly unitary and gives
\textbar V\textsubscript{ud}\textbar{} = 0.97467. That is 3.1\ensuremath{\sigma} above the
superallowed-decay value 0.97367 \ensuremath{\pm} 0.00032, so the form bets that the
first-row shortfall (the ``Cabibbo anomaly'') goes away. Belle II and
LHCb measurements of sin 2\ensuremath{\beta}, \ensuremath{\alpha} and \ensuremath{\gamma} at the 1\ensuremath{^\circ} level will decide it.

\subsection{1. Construction}

In the exact form of the Wolfenstein parametrisation (Wolfenstein 1983;
Charles et al.~2005), the inputs are \ensuremath{\lambda} = s\ensuremath{_{1}}\ensuremath{_{2}} = 1/\ensuremath{\surd}20 = 0.22361 and s\ensuremath{_{2}}\ensuremath{_{3}} =
A\ensuremath{\lambda}\ensuremath{^{2}} with A = \ensuremath{\surd}(2/3). They are followed by \ensuremath{\bar{\rho}} + i\ensuremath{\bar{\eta}} = R\textsubscript{u}
e\textsuperscript{i\ensuremath{\gamma}}, where R\textsubscript{u} = sin \ensuremath{\beta} because the
angle at \ensuremath{\alpha} is 90\ensuremath{^\circ}. The relation

s\ensuremath{_{1}}\ensuremath{_{3}} e\textsuperscript{i\ensuremath{\delta}} = (\ensuremath{\bar{\rho}} + i\ensuremath{\bar{\eta}}) s\ensuremath{_{1}}\ensuremath{_{2}} s\ensuremath{_{2}}\ensuremath{_{3}} c\ensuremath{_{2}}\ensuremath{_{3}} / {[}c\ensuremath{_{1}}\ensuremath{_{2}} (1 \ensuremath{-} (\ensuremath{\bar{\rho}} + i\ensuremath{\bar{\eta}})
s\ensuremath{_{2}}\ensuremath{_{3}}\ensuremath{^{2}}){]}

then gives an exactly unitary matrix in the standard form. Numerically,
the unitarity error is 2 \ensuremath{\times} 10\ensuremath{^{-}}\ensuremath{^{1}}\ensuremath{^{6}}.

\subsection{2. Comparison with measurement}

The measured values are PDG 2026 for
\textbar V\textsubscript{us}\textbar{} = 0.22431, \ensuremath{\alpha} = 84.1 (+3.7 \ensuremath{-}3.0)\ensuremath{^\circ},
\ensuremath{\gamma} = 66.4 (+2.7 \ensuremath{-}2.8)\ensuremath{^\circ} and the first row. For the other quantities we use
PDG 2024 (Navas et al.~2024): \textbar V\textsubscript{cb}\textbar{} =
0.0411 \ensuremath{\pm} 0.0012, \textbar V\textsubscript{ub}\textbar{} = (3.82 \ensuremath{\pm} 0.20)
\ensuremath{\times} 10\ensuremath{^{-}}\ensuremath{^{3}}, sin 2\ensuremath{\beta} = 0.709 \ensuremath{\pm} 0.011 and
\textbar V\textsubscript{td}/V\textsubscript{ts}\textbar{} = 0.207 \ensuremath{\pm}
0.0032 (0.001 statistical and 0.003 lattice, combined). We use direct
measurements only, not global-fit values, which already assume
unitarity.

{\def\LTcaptype{none} % do not increment counter
\begin{longtable}[]{@{}llll@{}}
\toprule\noalign{}
Quantity & Prediction & Measured & Pull \\
\midrule\noalign{}
\endhead
\bottomrule\noalign{}
\endlastfoot
\textbar V\textsubscript{us}\textbar{} & 0.22361 & 0.22431 \ensuremath{\pm} 0.00085 &
\ensuremath{-}0.83 \\
\textbar V\textsubscript{cb}\textbar{} & 0.04082 & 0.0411 \ensuremath{\pm} 0.0012 &
\ensuremath{-}0.23 \\
\textbar V\textsubscript{ub}\textbar{} & 0.00358 & 0.00382 \ensuremath{\pm} 0.00020 &
\ensuremath{-}1.19 \\
sin 2\ensuremath{\beta} & 0.70711 & 0.709 \ensuremath{\pm} 0.011 & \ensuremath{-}0.17 \\
\ensuremath{\alpha} (\ensuremath{^\circ}) & 90.0 & 84.1 (+3.7 \ensuremath{-}3.0) & +1.59 (upper error used) \\
\ensuremath{\gamma} (\ensuremath{^\circ}) & 67.5 & 66.4 (+2.7 \ensuremath{-}2.8) & +0.39 (2.8 used) \\
\textbar V\textsubscript{td}/V\textsubscript{ts}\textbar{} & 0.21035 &
0.207 \ensuremath{\pm} 0.0032 & +1.05 \\
\end{longtable}
}

\ensuremath{\chi}\ensuremath{^{2}} = 5.98. If nothing had been chosen this would correspond to p = 0.54,
but it has no such meaning here.

\subsection{3. What was chosen, and what fails}

The values 1/\ensuremath{\surd}20, \ensuremath{\surd}(2/3), 22.5\ensuremath{^\circ} and 67.5\ensuremath{^\circ} = 90\ensuremath{^\circ} \ensuremath{-} 22.5\ensuremath{^\circ} were selected
because they are simple. A 3 \ensuremath{\times} 3 down-quark mass matrix with fixed
weights built to produce this triangle gives \ensuremath{\chi}\ensuremath{^{2}} = 11.5 at \ensuremath{\mu} = 162.5 GeV
(\textbar V\textsubscript{us}\textbar{} +2.2\ensuremath{\sigma}) and 26 at
M\textsubscript{Z}, so the triangle is not yet tied to a mass matrix.
The form is offered as a target for a mass-matrix origin, not as one.

\subsection{4. Tests}

\begin{itemize}
\tightlist
\item
  sin 2\ensuremath{\beta} firmly away from 1/\ensuremath{\surd}2 = 0.7071.
\item
  \ensuremath{\alpha} firmly away from 90\ensuremath{^\circ}. Today the direct value is 1.6\ensuremath{\sigma} below.
\item
  \ensuremath{\gamma} firmly outside about 64--68\ensuremath{^\circ}.
\item
  A confirmed first-row shortfall, which would also rule out any model
  that keeps the CKM matrix exactly unitary at low energy.
\end{itemize}

\subsection{Data and code}

\texttt{analysis/quark\_triangle\_and\_press\_2026-11-25.py} in the TRPG
code archive, doi:10.5281/zenodo.24329239. This result is one component
of the framework in doi:10.5281/zenodo.24330763.

\subsection{Acknowledgements}

Due to medical necessity, AI language-model assistants were used, under
the authors' direction, to draft and edit text from the authors' ideas,
to carry out and check calculations, and to search the literature; every
result was checked by rerunning the published code. The authors reviewed
all content and take full responsibility for it.

\subsection{References}

\begin{itemize}
\tightlist
\item
  Antusch, S., King, S. F., Malinsky, M. and Spinrath, M. (2010). Quark
  mixing sum rules and the right unitarity triangle. \emph{Phys. Rev.~D}
  81, 033008.
\item
  Charles, J. et al.~(CKMfitter Group) (2005). CP violation and the CKM
  matrix: assessing the impact of the asymmetric B factories. \emph{Eur.
  Phys. J. C} 41, 1.
\item
  Jarlskog, C. (1985). Commutator of the quark mass matrices in the
  standard electroweak model and a measure of maximal CP
  nonconservation. \emph{Phys. Rev.~Lett.} 55, 1039.
\item
  Kobayashi, M. and Maskawa, T. (1973). CP-violation in the
  renormalizable theory of weak interaction. \emph{Prog. Theor. Phys.}
  49, 652.
\item
  Navas, S. et al.~(Particle Data Group) (2024). Review of particle
  physics. \emph{Phys. Rev.~D} 110, 030001; and the 2026 edition
  (summary tables and CKM review).
\item
  Wolfenstein, L. (1983). Parametrization of the Kobayashi--Maskawa
  matrix. \emph{Phys. Rev.~Lett.} 51, 1945.
\end{itemize}

\end{document}
